\begin{theorem}[Precongruence for restriction]
  \label{thm:precong-relabel}
  \lean{PLTS.WeakScheduler.relabel}\lean{PLTS.ProbabilisticForwardSimulation.relabel}\leanok
  \uses{def:forwardsim, def:relabel, def:weakstep}
  If $sys_C$ is forward-simulated by $sys_A$ along $R$, both over an extended
  alphabet $\mathrm{Label} \oplus \mathrm{Extra}$, then
  $\mathrm{relabel}(sys_C)$ is forward-simulated by $\mathrm{relabel}(sys_A)$
  along the same relation $R$.
\end{theorem}

\begin{proof}\leanok
  \uses{def:relabel, def:weakstep, def:weaktau}
  The initial condition is unchanged, restriction not touching initial states.
  For the step condition, a transition of $\mathrm{relabel}(sys_C)$ on $l$ is a
  transition of $sys_C$ on $\mathrm{inl}\,l$, which $R$ answers by a weak
  transition of $sys_A$.
  What has to be transported is the witness of that weak transition, since a
  weak transition is given by a scheduler whose emissions are typed by the
  label type, and the label type has changed.
  Transport it as follows.
  Ask the scheduler about the prefix read on the extended alphabet, and read its
  answer back along the map that sends $\mathrm{inl}\,l$ to $l$ and every
  auxiliary label to $\tau$.
  The transported scheduler is valid and emits $\tau$ alone, and it puts on
  each emission exactly the mass the original puts on the corresponding
  extended one: the auxiliary preimages of a base label all carry mass $0$,
  because a weak scheduler emits $\tau$ alone.
  Its halting mass is likewise the original's.
  So the two weak transitions have the same source and the same target
  distribution, and the coupling is the one $R$ already provides.
  The two silent labels name the same thing, $\mathrm{inl}\,\tau$ being the
  silent label of the extended alphabet, so the case split between a silent and
  a visible answer transports as it stands.
\end{proof}
